\IfFileExists{stacks-project.cls}{%
\documentclass{stacks-project}
}{%
\documentclass{amsart}
}

% For dealing with references we use the comment environment
\usepackage{verbatim}
\newenvironment{reference}{\comment}{\endcomment}
%\newenvironment{reference}{}{}
\newenvironment{slogan}{\comment}{\endcomment}
\newenvironment{history}{\comment}{\endcomment}

% For commutative diagrams we use Xy-pic
\usepackage[all]{xy}

% We use 2cell for 2-commutative diagrams.
\xyoption{2cell}
\UseAllTwocells

% We use multicol for the list of chapters between chapters
\usepackage{multicol}

% This is generally recommended for better output
\usepackage{lmodern}
\usepackage[T1]{fontenc}

% For cross-file-references
\usepackage{xr-hyper}

% Package for hypertext links:
\usepackage{hyperref}

% For any local file, say "hello.tex" you want to link to please
% use \externaldocument[hello-]{hello}
\externaldocument[introduction-]{introduction}
\externaldocument[conventions-]{conventions}
\externaldocument[sets-]{sets}
\externaldocument[categories-]{categories}
\externaldocument[topology-]{topology}
\externaldocument[sheaves-]{sheaves}
\externaldocument[sites-]{sites}
\externaldocument[stacks-]{stacks}
\externaldocument[fields-]{fields}
\externaldocument[algebra-]{algebra}
\externaldocument[brauer-]{brauer}
\externaldocument[homology-]{homology}
\externaldocument[derived-]{derived}
\externaldocument[simplicial-]{simplicial}
\externaldocument[more-algebra-]{more-algebra}
\externaldocument[smoothing-]{smoothing}
\externaldocument[modules-]{modules}
\externaldocument[sites-modules-]{sites-modules}
\externaldocument[injectives-]{injectives}
\externaldocument[cohomology-]{cohomology}
\externaldocument[sites-cohomology-]{sites-cohomology}
\externaldocument[dga-]{dga}
\externaldocument[dpa-]{dpa}
\externaldocument[sdga-]{sdga}
\externaldocument[hypercovering-]{hypercovering}
\externaldocument[schemes-]{schemes}
\externaldocument[constructions-]{constructions}
\externaldocument[properties-]{properties}
\externaldocument[morphisms-]{morphisms}
\externaldocument[coherent-]{coherent}
\externaldocument[divisors-]{divisors}
\externaldocument[limits-]{limits}
\externaldocument[varieties-]{varieties}
\externaldocument[topologies-]{topologies}
\externaldocument[descent-]{descent}
\externaldocument[perfect-]{perfect}
\externaldocument[more-morphisms-]{more-morphisms}
\externaldocument[flat-]{flat}
\externaldocument[groupoids-]{groupoids}
\externaldocument[more-groupoids-]{more-groupoids}
\externaldocument[etale-]{etale}
\externaldocument[chow-]{chow}
\externaldocument[intersection-]{intersection}
\externaldocument[pic-]{pic}
\externaldocument[weil-]{weil}
\externaldocument[adequate-]{adequate}
\externaldocument[dualizing-]{dualizing}
\externaldocument[duality-]{duality}
\externaldocument[discriminant-]{discriminant}
\externaldocument[derham-]{derham}
\externaldocument[local-cohomology-]{local-cohomology}
\externaldocument[algebraization-]{algebraization}
\externaldocument[curves-]{curves}
\externaldocument[resolve-]{resolve}
\externaldocument[models-]{models}
\externaldocument[functors-]{functors}
\externaldocument[equiv-]{equiv}
\externaldocument[pione-]{pione}
\externaldocument[etale-cohomology-]{etale-cohomology}
\externaldocument[proetale-]{proetale}
\externaldocument[relative-cycles-]{relative-cycles}
\externaldocument[more-etale-]{more-etale}
\externaldocument[trace-]{trace}
\externaldocument[crystalline-]{crystalline}
\externaldocument[spaces-]{spaces}
\externaldocument[spaces-properties-]{spaces-properties}
\externaldocument[spaces-morphisms-]{spaces-morphisms}
\externaldocument[decent-spaces-]{decent-spaces}
\externaldocument[spaces-cohomology-]{spaces-cohomology}
\externaldocument[spaces-limits-]{spaces-limits}
\externaldocument[spaces-divisors-]{spaces-divisors}
\externaldocument[spaces-over-fields-]{spaces-over-fields}
\externaldocument[spaces-topologies-]{spaces-topologies}
\externaldocument[spaces-descent-]{spaces-descent}
\externaldocument[spaces-perfect-]{spaces-perfect}
\externaldocument[spaces-more-morphisms-]{spaces-more-morphisms}
\externaldocument[spaces-flat-]{spaces-flat}
\externaldocument[spaces-groupoids-]{spaces-groupoids}
\externaldocument[spaces-more-groupoids-]{spaces-more-groupoids}
\externaldocument[bootstrap-]{bootstrap}
\externaldocument[spaces-pushouts-]{spaces-pushouts}
\externaldocument[spaces-chow-]{spaces-chow}
\externaldocument[groupoids-quotients-]{groupoids-quotients}
\externaldocument[spaces-more-cohomology-]{spaces-more-cohomology}
\externaldocument[spaces-simplicial-]{spaces-simplicial}
\externaldocument[spaces-duality-]{spaces-duality}
\externaldocument[formal-spaces-]{formal-spaces}
\externaldocument[restricted-]{restricted}
\externaldocument[spaces-resolve-]{spaces-resolve}
\externaldocument[formal-defos-]{formal-defos}
\externaldocument[defos-]{defos}
\externaldocument[cotangent-]{cotangent}
\externaldocument[examples-defos-]{examples-defos}
\externaldocument[algebraic-]{algebraic}
\externaldocument[examples-stacks-]{examples-stacks}
\externaldocument[stacks-sheaves-]{stacks-sheaves}
\externaldocument[criteria-]{criteria}
\externaldocument[artin-]{artin}
\externaldocument[quot-]{quot}
\externaldocument[stacks-properties-]{stacks-properties}
\externaldocument[stacks-morphisms-]{stacks-morphisms}
\externaldocument[stacks-limits-]{stacks-limits}
\externaldocument[stacks-cohomology-]{stacks-cohomology}
\externaldocument[stacks-perfect-]{stacks-perfect}
\externaldocument[stacks-introduction-]{stacks-introduction}
\externaldocument[stacks-more-morphisms-]{stacks-more-morphisms}
\externaldocument[stacks-geometry-]{stacks-geometry}
\externaldocument[moduli-]{moduli}
\externaldocument[moduli-curves-]{moduli-curves}
\externaldocument[examples-]{examples}
\externaldocument[exercises-]{exercises}
\externaldocument[guide-]{guide}
\externaldocument[desirables-]{desirables}
\externaldocument[coding-]{coding}
\externaldocument[obsolete-]{obsolete}
\externaldocument[fdl-]{fdl}
\externaldocument[index-]{index}

% Theorem environments.
%
\theoremstyle{plain}
\newtheorem{theorem}[subsection]{Theorem}
\newtheorem{proposition}[subsection]{Proposition}
\newtheorem{lemma}[subsection]{Lemma}

\theoremstyle{definition}
\newtheorem{definition}[subsection]{Definition}
\newtheorem{example}[subsection]{Example}
\newtheorem{exercise}[subsection]{Exercise}
\newtheorem{situation}[subsection]{Situation}

\theoremstyle{remark}
\newtheorem{remark}[subsection]{Remark}
\newtheorem{remarks}[subsection]{Remarks}

\numberwithin{equation}{subsection}

% Macros
%
\def\lim{\mathop{\mathrm{lim}}\nolimits}
\def\colim{\mathop{\mathrm{colim}}\nolimits}
\def\Spec{\mathop{\mathrm{Spec}}}
\def\Hom{\mathop{\mathrm{Hom}}\nolimits}
\def\Ext{\mathop{\mathrm{Ext}}\nolimits}
\def\SheafHom{\mathop{\mathcal{H}\!\mathit{om}}\nolimits}
\def\SheafExt{\mathop{\mathcal{E}\!\mathit{xt}}\nolimits}
\def\Sch{\mathit{Sch}}
\def\Mor{\mathop{\mathrm{Mor}}\nolimits}
\def\Ob{\mathop{\mathrm{Ob}}\nolimits}
\def\Sh{\mathop{\mathit{Sh}}\nolimits}
\def\NL{\mathop{N\!L}\nolimits}
\def\CH{\mathop{\mathrm{CH}}\nolimits}
\def\proetale{{pro\text{-}\acute{e}tale}}
\def\etale{{\acute{e}tale}}
\def\QCoh{\mathit{QCoh}}
\def\Ker{\mathop{\mathrm{Ker}}}
\def\Im{\mathop{\mathrm{Im}}}
\def\Coker{\mathop{\mathrm{Coker}}}
\def\Coim{\mathop{\mathrm{Coim}}}

% Boxtimes
%
\DeclareMathSymbol{\boxtimes}{\mathbin}{AMSa}{"02}

%
% Macros for moduli stacks/spaces
%
\def\QCohstack{\mathcal{QC}\!\mathit{oh}}
\def\Cohstack{\mathcal{C}\!\mathit{oh}}
\def\Spacesstack{\mathcal{S}\!\mathit{paces}}
\def\Quotfunctor{\mathrm{Quot}}
\def\Hilbfunctor{\mathrm{Hilb}}
\def\Curvesstack{\mathcal{C}\!\mathit{urves}}
\def\Polarizedstack{\mathcal{P}\!\mathit{olarized}}
\def\Complexesstack{\mathcal{C}\!\mathit{omplexes}}
% \Pic is the operator that assigns to X its picard group, usage \Pic(X)
% \Picardstack_{X/B} denotes the Picard stack of X over B
% \Picardfunctor_{X/B} denotes the Picard functor of X over B
\def\Pic{\mathop{\mathrm{Pic}}\nolimits}
\def\Picardstack{\mathcal{P}\!\mathit{ic}}
\def\Picardfunctor{\mathrm{Pic}}
\def\Deformationcategory{\mathcal{D}\!\mathit{ef}}
\usepackage{cleveref}
\crefname{originallemma}{original lemma}{original lemmas}
\Crefname{originallemma}{Original lemma}{Original lemmas}
\makeatletter
\newlabel{lemma-descend-opens}{{0827}{source}{}{sourcepointer-1}{}}
\newlabel{lemma-descend-opens@cref}{{[originallemma][0][]0827}{[1][0][]source}{}{}{}}
\makeatother
\crefname{originalreference}{original reference}{original references}
\Crefname{originalreference}{Original reference}{Original references}
\makeatletter
\newlabel{limits-lemma-finite-type-eventually-closed}{{0828}{source}{}{sourcepointer-2}{}}
\newlabel{limits-lemma-finite-type-eventually-closed@cref}{{[originalreference][0][]0828}{[1][0][]source}{}{}{}}
\makeatother
\crefname{originalreference}{original reference}{original references}
\Crefname{originalreference}{Original reference}{Original references}
\makeatletter
\newlabel{spaces-morphisms-lemma-diagonal-morphism-finite-type}{{084P}{source}{}{sourcepointer-3}{}}
\newlabel{spaces-morphisms-lemma-diagonal-morphism-finite-type@cref}{{[originalreference][0][]084P}{[1][0][]source}{}{}{}}
\makeatother
\crefname{originalreference}{original reference}{original references}
\Crefname{originalreference}{Original reference}{Original references}
\makeatletter
\newlabel{spaces-morphisms-lemma-fibre-product-after-map}{{03KN}{source}{}{sourcepointer-4}{}}
\newlabel{spaces-morphisms-lemma-fibre-product-after-map@cref}{{[originalreference][0][]03KN}{[1][0][]source}{}{}{}}
\makeatother
\crefname{originalreference}{original reference}{original references}
\Crefname{originalreference}{Original reference}{Original references}
\makeatletter
\newlabel{spaces-morphisms-lemma-finite-presentation-permanence}{{05WT}{source}{}{sourcepointer-5}{}}
\newlabel{spaces-morphisms-lemma-finite-presentation-permanence@cref}{{[originalreference][0][]05WT}{[1][0][]source}{}{}{}}
\makeatother
\crefname{originalreference}{original reference}{original references}
\Crefname{originalreference}{Original reference}{Original references}
\makeatletter
\newlabel{spaces-morphisms-lemma-permanence-finite-type}{{0462}{source}{}{sourcepointer-6}{}}
\newlabel{spaces-morphisms-lemma-permanence-finite-type@cref}{{[originalreference][0][]0462}{[1][0][]source}{}{}{}}
\makeatother
\title{Immersions and monomorphisms hold eventually over affine limit systems}
\author{The Stacks Project Authors}
\date{stacks-project:spaces-limits.tex:local-snapshot-20261012; GNU FDL 1.2 or later}
\begin{document}
\maketitle
\noindent\textit{Editorial note.} Counted span is the author's own proof environment, source lines 1475-1538, following the statement at 1451-1473.\par\medskip
\noindent\textit{Editorial note.} The '(details omitted)' in part (1) covers the choice of the affine etale cover U\_0 -> X\_0 x\_B W over the quasi-compact X\_0; flagged in the notes.\par\medskip
\noindent\textit{Editorial note.} Steps the author leaves to the reader are flagged here for the reviewer: source line(s) 1482 carry an "omitted"/"follows easily"/"immediate from the definitions" phrase. Those sentences are retained verbatim; the displayed argument around them stays inside the counted span.\par\medskip
\noindent\textit{Allowed standard prerequisites (editorial).}\begin{itemize}
\item Scheme-level eventual closedness. Limits, Lemma limits-lemma-finite-type-eventually-closed (cross-chapter, cited at \textasciitilde{}1487-1488). The scheme case consumed by part (1).
\item Descending opens along affine limit systems. lemma-descend-opens (same chapter, cited at \textasciitilde{}1497-1498). Part (2): writes U = lim U\_i with U\_i quasi-compact open in X\_i.
\item Finite presentation permanence. Morphisms of Spaces, Lemma spaces-morphisms-lemma-finite-presentation-permanence (cross-chapter, cited at \textasciitilde{}1511-1512). Part (3): U -> U\_0 is of finite presentation.
\item Diagonal and fibred product maps of spaces. Morphisms of Spaces, Lemmas spaces-morphisms-lemma-fibre-product-after-map, spaces-morphisms-lemma-permanence-finite-type, spaces-morphisms-lemma-diagonal-morphism-finite-type (cross-chapter, cited at \textasciitilde{}1522-1531). Part (4): transition maps of Z\_i are closed immersions; lfp of Y -> Z\_0.
\end{itemize}\medskip
\ifdefined\hypertarget\hypertarget{sourcepointer-1}{}\fi
\noindent\textit{Source reference.} Original lemma 0827: sources/stacks-spaces-limits/pinned-original-source.tex, source lines 1295--1295. Navigation pointer only: the counted proof cites this statement by \textbackslash{}ref as a standing input; the deduction the counted proof has to make is carried out inside the counted span, and no part of that deduction is delegated to this pointer.\par\medskip
\ifdefined\hypertarget\hypertarget{sourcepointer-2}{}\fi
\noindent\textit{Source reference.} Original reference 0828: focused-seeds/geometry/S9-spaces-limits-finite-type-eventually-closed/ref-tags.txt, source lines 5--5. Navigation pointer only: the counted proof cites this statement by \textbackslash{}ref as a standing input; the deduction the counted proof has to make is carried out inside the counted span, and no part of that deduction is delegated to this pointer.\par\medskip
\ifdefined\hypertarget\hypertarget{sourcepointer-3}{}\fi
\noindent\textit{Source reference.} Original reference 084P: focused-seeds/geometry/S9-spaces-limits-finite-type-eventually-closed/ref-tags.txt, source lines 6--6. Navigation pointer only: the counted proof cites this statement by \textbackslash{}ref as a standing input; the deduction the counted proof has to make is carried out inside the counted span, and no part of that deduction is delegated to this pointer.\par\medskip
\ifdefined\hypertarget\hypertarget{sourcepointer-4}{}\fi
\noindent\textit{Source reference.} Original reference 03KN: focused-seeds/geometry/S9-spaces-limits-finite-type-eventually-closed/ref-tags.txt, source lines 7--7. Navigation pointer only: the counted proof cites this statement by \textbackslash{}ref as a standing input; the deduction the counted proof has to make is carried out inside the counted span, and no part of that deduction is delegated to this pointer.\par\medskip
\ifdefined\hypertarget\hypertarget{sourcepointer-5}{}\fi
\noindent\textit{Source reference.} Original reference 05WT: focused-seeds/geometry/S9-spaces-limits-finite-type-eventually-closed/ref-tags.txt, source lines 8--8. Navigation pointer only: the counted proof cites this statement by \textbackslash{}ref as a standing input; the deduction the counted proof has to make is carried out inside the counted span, and no part of that deduction is delegated to this pointer.\par\medskip
\ifdefined\hypertarget\hypertarget{sourcepointer-6}{}\fi
\noindent\textit{Source reference.} Original reference 0462: focused-seeds/geometry/S9-spaces-limits-finite-type-eventually-closed/ref-tags.txt, source lines 9--9. Navigation pointer only: the counted proof cites this statement by \textbackslash{}ref as a standing input; the deduction the counted proof has to make is carried out inside the counted span, and no part of that deduction is delegated to this pointer.\par\medskip
\begin{lemma}
\label{lemma-finite-type-eventually-closed}
Let $S$ be a scheme. Let $B$ be an algebraic space over $S$.
Let $X = \lim X_i$ be a directed limit of
algebraic spaces over $B$ with affine transition morphisms.
Let $Y \to X$ be a morphism of algebraic spaces over $B$.
\begin{enumerate}
\item If $Y \to X$ is a closed immersion, $X_i$ quasi-compact, and
$Y \to B$ locally of finite type, then $Y \to X_i$ is a closed immersion
for $i$ large enough.
\item If $Y \to X$ is an immersion, $X_i$ quasi-separated, $Y \to B$
locally of finite type, and $Y$ quasi-compact, then $Y \to X_i$ is an
immersion for $i$ large enough.
\item If $Y \to X$ is an isomorphism, $X_i$ quasi-compact,
$X_i \to B$ locally of finite type, the transition morphisms
$X_{i'} \to X_i$ are closed immersions, and $Y \to B$ is locally
of finite presentation, then $Y \to X_i$ is an isomorphism for $i$
large enough.
\item If $Y \to X$ is a monomorphism, $X_i$ quasi-separated,
$Y \to B$ locally of finite type, and $Y$ quasi-compact, then
$Y \to X_i$ is a monomorphism for $i$ large enough.
\end{enumerate}
\end{lemma}

\begin{proof}
Proof of (1). Choose $0 \in I$. As $X_0$ is quasi-compact, we can choose
an affine scheme $W$ and an \'etale morphism $W \to B$ such that the image of
$|X_0| \to |B|$ is contained in $|W| \to |B|$. 
Choose an affine scheme $U_0$ and an \'etale morphism
$U_0 \to X_0 \times_B W$ such that $U_0 \to X_0$ is surjective.
(This is possible by our choice of $W$ and the fact that $X_0$ is
quasi-compact; details omitted.)
Let $V \to Y$, resp.\ $U \to X$, resp.\ $U_i \to X_i$ be the base change
of $U_0 \to X_0$ (for $i \geq 0$). It suffices to prove that $V \to U_i$
is a closed immersion for $i$ sufficiently large. Thus we reduce
to proving the result for $V \to U = \lim U_i$ over $W$. This follows from
the case of schemes, which is
Limits, Lemma \ref{limits-lemma-finite-type-eventually-closed}.

\medskip\noindent
Proof of (2). Choose $0 \in I$. Choose a quasi-compact open subspace
$X'_0 \subset X_0$ such that $Y \to X_0$ factors through $X'_0$.
After replacing $X_i$ by the inverse image of $X'_0$ for $i \geq 0$
we may assume all $X_i'$ are quasi-compact and quasi-separated.
Let $U \subset X$ be a quasi-compact open such that $Y \to X$ factors
through a closed immersion $Y \to U$ ($U$ exists as $Y$ is quasi-compact). By
Lemma \ref{lemma-descend-opens}
we may assume that $U = \lim U_i$ with $U_i \subset X_i$ quasi-compact
open. By part (1) we see that $Y \to U_i$ is a closed immersion for some
$i$. Thus (2) holds.

\medskip\noindent
Proof of (3). Choose $0 \in I$. Choose an affine scheme $U_0$
and a surjective \'etale morphism $U_0 \to X_0$.
Set $U_i = X_i \times_{X_0} U_0$,
$U = X \times_{X_0} U_0 = Y \times_{X_0} U_0$. Then $U = \lim U_i$ is a
limit of affine schemes, the transition maps of the system are closed
immersions, and $U \to U_0$ is of finite presentation (because
$U \to B$ is locally of finite presentation and $U_0 \to B$ is locally
of finite type and
Morphisms of Spaces, Lemma
\ref{spaces-morphisms-lemma-finite-presentation-permanence}).
Thus we've reduced to the following algebra fact: If $A = \lim A_i$
is a directed colimit of $R$-algebras with surjective transition
maps and $A$ of finite presentation over $A_0$, then $A = A_i$ for
some $i$. Namely, write $A = A_0/(f_1, \ldots, f_n)$. Pick $i$ such
that $f_1, \ldots, f_n$ map to zero under the surjective map $A_0 \to A_i$.

\medskip\noindent
Proof of (4). Set $Z_i = Y \times_{X_i} Y$.
As the transition morphisms $X_{i'} \to X_i$ are affine hence separated,
the transition morphisms $Z_{i'} \to Z_i$ are closed immersions, see
Morphisms of Spaces, Lemma
\ref{spaces-morphisms-lemma-fibre-product-after-map}.
We have $\lim Z_i = Y \times_X Y = Y$ as $Y \to X$ is a monomorphism.
Choose $0 \in I$. Since $Y \to X_0$ is locally of finite type
(Morphisms of Spaces, Lemma
\ref{spaces-morphisms-lemma-permanence-finite-type})
the morphism $Y \to Z_0$ is locally of finite presentation
(Morphisms of Spaces, Lemma
\ref{spaces-morphisms-lemma-diagonal-morphism-finite-type}).
The morphisms $Z_i \to Z_0$ are locally of finite type
(they are closed immersions).
Finally, $Z_i = Y \times_{X_i} Y$ is quasi-compact as
$X_i$ is quasi-separated and $Y$ is quasi-compact.
Thus part (3) applies to $Y = \lim_{i \geq 0} Z_i$ over $Z_0$
and we conclude $Y = Z_i$ for some $i$. This proves (4) and the lemma.
\end{proof}

\end{document}
